\documentclass[11pt]{article}

\usepackage[a4paper,margin=1in]{geometry}
\usepackage{amsmath,amssymb,booktabs,longtable,array}
\usepackage{hyperref}
\usepackage[T1]{fontenc}
\usepackage[utf8]{inputenc}

\title{GENESIS Conditional UNet: Detailed Architecture and Conditioning Notes}
\author{}
\date{}

\begin{document}
\maketitle
\tableofcontents
\newpage

\section{Purpose of this document}
This note is a code-aligned description of the conditional UNet used in GENESIS for
flow matching and diffusion-style training. The goal is to document \emph{exactly}
how a single input tensor flows through the network, where conditioning enters,
what each configurable option does, and how the normalization / attention / residual
blocks are implemented.

The description follows the implementation in \texttt{models/unet.py},
\texttt{models/embeddings.py}, \texttt{flow\_matching/flows.py}, and
\texttt{flow\_matching/samplers.py}.

\section{High-level overview}
The model is a ResNet-style UNet for 3-channel cosmology maps of size
\[
x \in \mathbb{R}^{3 \times 256 \times 256}.
\]
It predicts a vector field
\[
v(x_t, t, c) \in \mathbb{R}^{3 \times 256 \times 256},
\]
where:
\begin{itemize}
  \item $t \in [0,1]$ is the flow time.
  \item $c \in \mathbb{R}^6$ is the cosmology conditioning vector.
  \item $x_t$ is the intermediate state along the flow / interpolation path.
\end{itemize}

The forward path is:
\[
x \to \text{input stem} \to \text{encoder} \to \text{bottleneck} \to \text{decoder} \to \text{output head}.
\]

The conditioning path is:
\[
t \to t_{\mathrm{emb}}, \qquad c \to c_{\mathrm{emb}}, \qquad
(t_{\mathrm{emb}}, c_{\mathrm{emb}}) \to \text{joint embedding}.
\]

\section{Constructor options}
The UNet constructor exposes the following main options.

\begin{longtable}{>{\raggedright\arraybackslash}p{0.20\textwidth}
                  >{\raggedright\arraybackslash}p{0.22\textwidth}
                  >{\raggedright\arraybackslash}p{0.18\textwidth}
                  >{\raggedright\arraybackslash}p{0.34\textwidth}}
\toprule
\textbf{Option} & \textbf{Default / values} & \textbf{Type} & \textbf{Meaning} \\
\midrule
\endhead
\texttt{preset} & \texttt{S}, \texttt{B}, \texttt{L} & string & Overrides the width / attention / block depth preset. \\
\texttt{img\_size} & 256 & int & Input spatial resolution. Used to build the stage bookkeeping. \\
\texttt{in\_channels} & 3 & int & Number of input map channels. \\
\texttt{cond\_dim} & 6 & int & Dimension of the raw conditioning vector. \\
\texttt{base\_channels} & 128 (if no preset) & int & Base width of the UNet. \\
\texttt{channel\_mult} & \texttt{(1,2,4,4)} if no preset; preset B uses \texttt{[1,2,3,4]} & list[int] & Channel multipliers per encoder / decoder stage. \\
\texttt{num\_res\_blocks} & 2 (preset B) & int & Number of ResBlocks per encoder stage. Decoder has one extra ResBlock per stage. \\
\texttt{num\_heads} & 8 (preset B) & int & Number of attention heads. Must divide the channel width of attention stages. \\
\texttt{attention\_resolution} & 32 & int & Attention is inserted only when the bookkeeping resolution is at or below this value. \\
\texttt{channel\_se} & True & bool & Enables bottleneck ChannelSE. \\
\texttt{dropout} & 0.0 & float & Dropout probability inside ResBlocks. \\
\texttt{circular\_conv} & False & bool & Uses circular padding / periodic convolution. \\
\texttt{upsample\_mode} & \texttt{nearest} & string & Decoder upsampling mode: \texttt{nearest} or \texttt{bilinear}. \\
\texttt{cond\_context\_mode} & \texttt{cond} & string & Auxiliary conditioning context for cross-attention / per-scale projection: \texttt{cond} or \texttt{joint}. \\
\texttt{cross\_attn\_cond} & True & bool & Enables cross-attention blocks at attention-enabled resolutions. \\
\texttt{per\_scale\_cond} & True & bool & Enables stage-specific condition projections. \\
\texttt{cond\_depth} & 4 & int & Depth of the condition MLP. \\
\bottomrule
\end{longtable}

\subsection{Presets}
The presets are defined as:
\begin{itemize}
  \item \texttt{S}: \texttt{base\_channels=64}, \texttt{channel\_mult=[1,2,3,4]}, \texttt{num\_heads=4}, \texttt{num\_res\_blocks=2}
  \item \texttt{B}: \texttt{base\_channels=128}, \texttt{channel\_mult=[1,2,3,4]}, \texttt{num\_heads=8}, \texttt{num\_res\_blocks=2}
  \item \texttt{L}: \texttt{base\_channels=256}, \texttt{channel\_mult=[1,2,3,4]}, \texttt{num\_heads=16}, \texttt{num\_res\_blocks=3}
\end{itemize}
If a preset is provided, it overrides the corresponding explicit constructor defaults.

\section{Embeddings}
\subsection{Timestep embedding}
The scalar time $t$ is mapped to a sinusoidal embedding and then passed through a two-layer MLP:
\[
t \mapsto t_{\mathrm{emb}} \in \mathbb{R}^{D}, \qquad D = 4 \times \texttt{base\_channels}.
\]
The sinusoidal sub-embedding uses dimension 256 in the implementation, and the MLP output dimension is the main embedding width.

\subsection{Condition embedding}
The 6-dimensional cosmology parameter vector is mapped through a deeper MLP:
\[
c \mapsto c_{\mathrm{emb}} \in \mathbb{R}^{D}.
\]
The default depth is four linear layers (\texttt{cond\_depth=4}) with SiLU activations between hidden layers.

\subsection{Joint embedding}
The model does \emph{not} simply add the time and condition embeddings.
Instead, it concatenates them and projects them with an MLP:
\[
\mathrm{joint} = \mathrm{JointEmbedding}(t_{\mathrm{emb}}, c_{\mathrm{emb}}) \in \mathbb{R}^{D}.
\]
This preserves the two signals separately until the final projection and reduces information loss.

\subsection{Alternative token embedding}
The codebase also contains \texttt{ConditionTokenEmbedding}, which can convert the scalar condition vector into a small sequence of tokens. This module is currently not used by the main UNet forward pass, but it is useful for token-based conditioning experiments.

\section{Normalization layers}
\subsection{GroupNorm}
The basic normalization primitive is:
\[
\mathrm{GroupNorm}(32, C, \epsilon=10^{-6}).
\]
This means:
\begin{itemize}
  \item the channel dimension is split into 32 groups,
  \item each group is normalized independently per sample,
  \item the operation does \emph{not} mix samples in the batch,
  \item affine parameters $\gamma_c$ and $\beta_c$ are learnable by default (PyTorch default behavior).
\end{itemize}

For an input tensor $x \in \mathbb{R}^{B \times C \times H \times W}$, GroupNorm normalizes each channel group by the mean and variance computed over channels in the group and over spatial positions $(H,W)$, for each sample separately. The small epsilon stabilizes the division.

\subsection{LayerNorm in cross-attention}
The conditioning vector used by cross-attention is normalized with \texttt{LayerNorm}. This is applied to the context vector, not to the spatial feature map.

\section{Residual block with AdaGN}
The core building block is the ResBlock:
\[
\text{ResBlock}(x, e).
\]
It uses AdaGN-style conditioning:
\[
\tilde{x} = \mathrm{GN}(x), \qquad
(\gamma, \beta) = f_{\mathrm{emb}}(e),
\qquad
\mathrm{AdaGN}(x, e) = \tilde{x} \odot (1+\gamma) + \beta.
\]

\subsection{Internal structure of a ResBlock}
For input feature map $x$ and embedding $e$:
\begin{enumerate}
  \item Normalize and activate:
  \[
  h \leftarrow \mathrm{SiLU}(\mathrm{GN}_1(x)).
  \]
  \item Apply a $3 \times 3$ convolution:
  \[
  h \leftarrow \mathrm{Conv}_{3 \times 3}(h).
  \]
  \item Project the embedding:
  \[
  (\gamma, \beta) = \mathrm{Linear}(\mathrm{SiLU}(e)).
  \]
  The linear layer outputs $2 \cdot C_{\mathrm{out}}$ values, which are split into scale and shift.
  \item Apply AdaGN:
  \[
  h \leftarrow \mathrm{GN}_2(h) \odot (1+\gamma) + \beta.
  \]
  \item Apply dropout and SiLU:
  \[
  h \leftarrow \mathrm{Dropout}(\mathrm{SiLU}(h)).
  \]
  \item Apply a second $3 \times 3$ convolution.
  \item Add the residual skip connection:
  \[
  y = h + \mathrm{Skip}(x).
  \]
\end{enumerate}

\subsection{Residual skip path}
If $C_{\mathrm{in}} \neq C_{\mathrm{out}}$, the skip path is a $1 \times 1$ convolution.
Otherwise it is the identity.

\subsection{Initialization}
The second convolution in each ResBlock is zero-initialized, so the block starts close to an identity mapping.
The output convolution is also zero-initialized. This helps stabilize early training.

\section{Self-attention block}
The self-attention block is applied to spatial feature maps. It uses:
\begin{itemize}
  \item \texttt{GroupNorm(32, C, eps=1e-6)} on the feature map,
  \item a $1 \times 1$ convolution to produce $Q,K,V$,
  \item PyTorch's \texttt{scaled\_dot\_product\_attention},
  \item a final $1 \times 1$ output projection.
\end{itemize}

\subsection{Shape logic}
For $x \in \mathbb{R}^{B \times C \times H \times W}$:
\[
Q,K,V \in \mathbb{R}^{B \times C \times H \times W}
\]
are reshaped into head format:
\[
\mathbb{R}^{B \times N_h \times (HW) \times d_h},
\qquad d_h = \frac{C}{N_h}.
\]
The attention is then computed across spatial positions. The output is reshaped back to
\[
\mathbb{R}^{B \times C \times H \times W}
\]
and added residually to the input.

\section{Cross-attention block}
Cross-attention injects the conditioning vector directly into a spatial feature map.

\subsection{Roles of query, key, and value}
\begin{itemize}
  \item Query $Q$: comes from the spatial feature map.
  \item Key $K$: comes from the conditioning context.
  \item Value $V$: comes from the conditioning context.
\end{itemize}

The context is either:
\[
\text{context} = c_{\mathrm{emb}}
\quad \text{or} \quad
\text{context} = \mathrm{joint},
\]
depending on \texttt{cond\_context\_mode}.

\subsection{Implementation details}
For a feature map $x$:
\begin{enumerate}
  \item Normalize the feature map with GroupNorm.
  \item Normalize the context with LayerNorm.
  \item Project $x$ to queries with a $1 \times 1$ convolution.
  \item Project the context to keys and values with linear layers.
  \item Run multi-head scaled dot-product attention.
  \item Apply a $1 \times 1$ output projection and add the residual.
\end{enumerate}

The output projection is zero-initialized, so the block starts as a near-identity module.

\section{Channel Squeeze-and-Excitation}
The ChannelSE module is a bottleneck-only channel recalibration block. It does \emph{not} take the condition vector as an input.

\subsection{Mechanism}
Given a bottleneck feature map $x \in \mathbb{R}^{B \times C \times H \times W}$:
\begin{enumerate}
  \item Global average pool each channel:
  \[
  x \mapsto \mathrm{GAP}(x) \in \mathbb{R}^{B \times C}.
  \]
  \item Pass through a small MLP:
  \[
  C \to \max(C/r, 4) \to C,
  \]
  where $r=16$ by default.
  \item Apply sigmoid gating:
  \[
  g \in (0,1)^C.
  \]
  \item Reweight the channels:
  \[
  y = x \odot g.
  \]
\end{enumerate}

\subsection{Interpretation}
ChannelSE does not move spatial information. It only says which feature channels should be emphasized or suppressed after the bottleneck has already compressed the spatial structure.

\section{Up-sampling and down-sampling}
\subsection{Down-sampling}
Down-sampling is implemented as a $3 \times 3$ convolution with stride 2:
\[
H \times W \to \frac{H}{2} \times \frac{W}{2}.
\]
The channel count stays the same in the down-sample layer itself.

\subsection{Up-sampling}
Up-sampling uses:
\begin{itemize}
  \item \texttt{nearest}: nearest-neighbor interpolation by factor 2,
  \item \texttt{bilinear}: bilinear interpolation by factor 2 with \texttt{align\_corners=False},
\end{itemize}
followed by a $3 \times 3$ convolution.

\subsection{Interpretation}
The interpolation step only enlarges the spatial grid. The following convolution then reworks the feature content.
This separation makes the decoder behavior easier to control and stabilizes training.

\section{Encoder forward pass in detail}
The encoder starts from:
\[
h_0 = \mathrm{InputConv}(x) \in \mathbb{R}^{128 \times 256 \times 256}
\]
for preset B.

The stage widths under preset B are:
\[
128,\;256,\;384,\;512.
\]

\subsection{Stage 0: $256 \times 256$}
\begin{itemize}
  \item Input width: 128.
  \item Output width: 128.
  \item Two ResBlocks.
  \item No attention if threshold is 32 or 64, because this stage is at full resolution.
  \item Downsample to $128 \times 128$.
\end{itemize}

\subsection{Stage 1: $128 \times 128$}
\begin{itemize}
  \item The first ResBlock expands channels from 128 to 256.
  \item The second ResBlock keeps 256 channels.
  \item No attention under thresholds 32 or 64.
  \item Downsample to $64 \times 64$.
\end{itemize}

\subsection{Stage 2: $64 \times 64$}
\begin{itemize}
  \item The first ResBlock expands channels from 256 to 384.
  \item The second ResBlock keeps 384 channels.
  \item Attention appears here only when the threshold is at least 64.
  \item Downsample to $32 \times 32$.
\end{itemize}

\subsection{Stage 3: $32 \times 32$}
\begin{itemize}
  \item The first ResBlock expands channels from 384 to 512.
  \item The second ResBlock keeps 512 channels.
  \item Attention appears here for thresholds 32 and 64.
  \item Downsample to the bottleneck at $16 \times 16$.
\end{itemize}

\section{Bottleneck in detail}
The bottleneck is the most compressed feature representation:
\[
512 \times 16 \times 16.
\]

The sequence is:
\begin{enumerate}
  \item ResBlock with bottleneck embedding.
  \item Self-attention.
  \item Cross-attention, if enabled.
  \item ChannelSE, if enabled.
  \item Second ResBlock with the same bottleneck embedding.
\end{enumerate}

This is the global reasoning center of the UNet. The spatial grid is small enough that attention can capture broad dependencies efficiently.

\section{Decoder forward pass in detail}
The decoder mirrors the encoder in reverse order. For each decoder stage:
\begin{enumerate}
  \item Upsample the current feature map by factor 2.
  \item Concatenate the corresponding encoder skip connection along channels.
  \item Apply \texttt{num\_res\_blocks + 1} ResBlocks.
  \item Apply attention blocks at the stages whose bookkeeping resolution is below the attention threshold.
\end{enumerate}

For preset B, \texttt{num\_res\_blocks = 2}, so each decoder stage has 3 ResBlocks.

\subsection{Decoder stage 0: $16 \times 16 \to 32 \times 32$}
\begin{itemize}
  \item Upsample: $512 \times 16 \times 16 \to 512 \times 32 \times 32$.
  \item Skip: encoder stage 3 output, also $512 \times 32 \times 32$.
  \item Concatenate to $1024 \times 32 \times 32$.
  \item First ResBlock reduces width back to 512.
  \item Remaining ResBlocks refine the 512-channel representation.
  \item Attention is active at this stage for thresholds 32 and above.
\end{itemize}

\subsection{Decoder stage 1: $32 \times 32 \to 64 \times 64$}
\begin{itemize}
  \item Upsample to $512 \times 64 \times 64$.
  \item Skip: encoder stage 2 output, $384 \times 64 \times 64$.
  \item Concatenate to $896 \times 64 \times 64$.
  \item First ResBlock reduces width back to 384.
  \item Remaining ResBlocks refine the 384-channel representation.
  \item Attention is active already for threshold 32, because this stage is inserted when the bookkeeping resolution has reached 32.
\end{itemize}

\subsection{Decoder stage 2: $64 \times 64 \to 128 \times 128$}
\begin{itemize}
  \item Upsample to $384 \times 128 \times 128$.
  \item Skip: encoder stage 1 output, $256 \times 128 \times 128$.
  \item Concatenate to $640 \times 128 \times 128$.
  \item First ResBlock reduces width back to 256.
  \item Remaining ResBlocks keep 256 channels.
  \item Attention is active only when the threshold reaches 64 or higher.
\end{itemize}

\subsection{Decoder stage 3: $128 \times 128 \to 256 \times 256$}
\begin{itemize}
  \item Upsample to $256 \times 256 \times 256$.
  \item Skip: encoder stage 0 output, $128 \times 256 \times 256$.
  \item Concatenate to $384 \times 256 \times 256$.
  \item First ResBlock reduces width back to 128.
  \item Remaining ResBlocks keep 128 channels.
  \item No attention at this resolution.
\end{itemize}

\section{Output head}
After the decoder, the tensor has shape
\[
128 \times 256 \times 256.
\]
The output head is:
\[
\mathrm{out} = \mathrm{Conv}_{3 \times 3}(\mathrm{SiLU}(\mathrm{GN}(h))).
\]
This returns the final 3-channel tensor:
\[
3 \times 256 \times 256.
\]

The output convolution is zero-initialized, again to make the network start close to a simple baseline mapping.

\section{How conditioning is routed}
There are three distinct conditioning routes.

\subsection{Route 1: Root joint embedding}
The base conditioning signal is:
\[
\mathrm{joint} = \mathrm{JointEmbedding}(t_{\mathrm{emb}}, c_{\mathrm{emb}}).
\]
This always exists and always feeds the ResBlocks through AdaGN.

\subsection{Route 2: Stage-specific projections}
If \texttt{per\_scale\_cond = true}, each stage uses:
\[
\ell_i = \mathrm{joint} + W_i(\text{context}),
\]
where:
\begin{itemize}
  \item \texttt{context = c\_emb} if \texttt{cond\_context\_mode = cond},
  \item \texttt{context = joint} if \texttt{cond\_context\_mode = joint}.
\end{itemize}

This is the stage-specific adapter path.

\subsection{Route 3: Cross-attention}
If \texttt{cross\_attn\_cond = true}, attention-enabled blocks can additionally query the same context vector directly.

\subsection{Condition-context mode}
The option \texttt{cond\_context\_mode} controls \emph{which vector is used as auxiliary context} in routes 2 and 3:
\begin{itemize}
  \item \texttt{cond}: use the pure condition embedding $c_{\mathrm{emb}}$.
  \item \texttt{joint}: use the combined time-condition embedding \texttt{joint}.
\end{itemize}

This option does \emph{not} remove the root joint embedding. It only changes the auxiliary context path.

\section{Attention placement logic}
Attention blocks are inserted according to a bookkeeping resolution variable \texttt{res}. The model adds attention when:
\[
\texttt{res} \leq \texttt{attention\_resolution}.
\]
This means:
\begin{itemize}
  \item higher attention threshold \(\Rightarrow\) attention starts earlier at larger spatial sizes,
  \item lower threshold \(\Rightarrow\) attention is restricted to smaller spatial sizes.
\end{itemize}

This keeps the computation manageable while still allowing the model to reason globally at low resolution.

\section{Flow matching usage}
The UNet is trained as the velocity / vector-field predictor inside a flow matching objective.

\subsection{OT flow matching}
The default OT-style flow matching objective samples:
\[
t \sim \mathcal{U}(0,1), \qquad x_1 \sim \mathcal{N}(0, I),
\]
and forms
\[
x_t = (1-t)x_0 + t x_1.
\]
The target vector field is:
\[
v^\star = x_1 - x_0.
\]
The model is trained to minimize:
\[
\mathcal{L} = \mathbb{E}\big[\|v_\theta(x_t, t, c) - (x_1 - x_0)\|^2\big].
\]

\subsection{Classifier-free guidance during training}
With probability \texttt{cfg\_prob}, the conditioning vector is zeroed before being passed to the model. This teaches the model both conditional and unconditional behavior.

\subsection{Sampling}
At inference, the samplers integrate the learned vector field from noise to data, typically from $t=1$ down to $t=0$. If classifier-free guidance is used, the sampler evaluates both:
\[
v_{\mathrm{cond}} = f_\theta(x,t,c), \qquad v_{\mathrm{uncond}} = f_\theta(x,t,0),
\]
and combines them with a guidance scale.

\section{What each option actually changes}
\begin{itemize}
  \item \texttt{preset}: changes the width and attention budget of the whole network.
  \item \texttt{attention\_resolution}: changes where attention begins.
  \item \texttt{circular\_conv}: changes the boundary condition of convolutions.
  \item \texttt{upsample\_mode}: changes how decoder resolution is increased.
  \item \texttt{cond\_context\_mode}: changes whether auxiliary conditioning uses $c_{\mathrm{emb}}$ or \texttt{joint}.
  \item \texttt{cross\_attn\_cond}: toggles explicit condition injection via cross-attention.
  \item \texttt{per\_scale\_cond}: toggles stage-wise adapters for the condition.
  \item \texttt{channel\_se}: toggles bottleneck channel reweighting.
  \item \texttt{cond\_depth}: changes the depth of the condition MLP.
  \item \texttt{num\_heads}: changes the number of attention heads.
  \item \texttt{dropout}: changes the internal ResBlock dropout.
\end{itemize}

\section{Short conceptual summary}
The model can be summarized as follows:
\begin{enumerate}
  \item Encode time and cosmology parameters separately.
  \item Fuse them into a root joint embedding.
  \item Inject that signal into every residual block via AdaGN.
  \item Optionally inject the condition again through cross-attention.
  \item Optionally modulate each spatial scale with its own projection.
  \item Compress the map to a small bottleneck, apply attention / SE, and decode with skip connections.
  \item Predict a vector field with the same spatial shape as the input map.
\end{enumerate}

\section{Worked forward walkthrough: preset B, full conditioning}
This section traces the exact forward pass in the style of a code execution log.
The assumed configuration is:
\begin{itemize}
  \item \texttt{preset = B}
  \item \texttt{attention\_resolution = 32}
  \item \texttt{circular\_conv = true}
  \item \texttt{upsample\_mode = nearest}
  \item \texttt{cond\_context\_mode = cond}
  \item \texttt{cross\_attn\_cond = true}
  \item \texttt{per\_scale\_cond = true}
  \item \texttt{channel\_se = true}
  \item \texttt{dropout = 0.0}
  \item \texttt{cond\_depth = 4}
\end{itemize}

\subsection{Input and embeddings}
The model receives:
\[
x \in \mathbb{R}^{3 \times 256 \times 256}, \qquad t \in [0,1], \qquad c \in \mathbb{R}^6.
\]
The embeddings are created in the following order:
\begin{enumerate}
  \item \texttt{t\_emb = TimestepEmbedding(t)}.
  \item \texttt{c\_emb = ConditionEmbedding(c)}.
  \item \texttt{joint = JointEmbedding(t\_emb, c\_emb)}.
\end{enumerate}
For this walkthrough the auxiliary context used by cross-attention is
\[
\text{cross\_ctx} = c_{\mathrm{emb}},
\]
because \texttt{cond\_context\_mode = cond}.

\subsection{Stem}
The stem convolution maps
\[
3 \times 256 \times 256 \to 128 \times 256 \times 256.
\]
This is a pure feature lift. No attention is used at the stem.

\subsection{Encoder stage 0: $256 \times 256$}
The first encoder stage operates at full resolution.
\paragraph{Conditioning.}
Because \texttt{per\_scale\_cond = true}, the stage-specific embedding is
\[
\ell_0 = \mathrm{joint} + W_0(c_{\mathrm{emb}}).
\]
\paragraph{Block 1.}
The first ResBlock keeps the width:
\[
128 \times 256 \times 256 \to 128 \times 256 \times 256.
\]
Inside the block:
\begin{itemize}
  \item \texttt{GroupNorm(32, 128, eps=1e-6)} normalizes channels in 32 groups.
  \item SiLU introduces nonlinearity.
  \item A $3 \times 3$ convolution mixes local spatial context.
  \item \texttt{emb\_proj} converts $\ell_0$ into scale and shift.
  \item AdaGN applies $\mathrm{GN}(h)\odot (1+\gamma) + \beta$.
  \item A second convolution and a residual skip finish the block.
\end{itemize}
At this high resolution, attention is not yet active.

\paragraph{Block 2.}
The second ResBlock also keeps the width:
\[
128 \times 256 \times 256 \to 128 \times 256 \times 256.
\]

\paragraph{Downsample.}
The stage ends with
\[
128 \times 256 \times 256 \to 128 \times 128 \times 128.
\]

\subsection{Encoder stage 1: $128 \times 128$}
The next stage receives
\[
128 \times 128 \times 128,
\]
with embedding
\[
\ell_1 = \mathrm{joint} + W_1(c_{\mathrm{emb}}).
\]
\paragraph{Block 1.}
The first ResBlock expands channels:
\[
128 \times 128 \times 128 \to 256 \times 128 \times 128.
\]
\paragraph{Block 2.}
The second ResBlock keeps the width:
\[
256 \times 128 \times 128 \to 256 \times 128 \times 128.
\]
\paragraph{Downsample.}
The stage ends with
\[
256 \times 128 \times 128 \to 256 \times 64 \times 64.
\]

\subsection{Encoder stage 2: $64 \times 64$}
The third stage receives
\[
256 \times 64 \times 64,
\]
with embedding
\[
\ell_2 = \mathrm{joint} + W_2(c_{\mathrm{emb}}).
\]
\paragraph{Block 1.}
The first ResBlock expands channels:
\[
256 \times 64 \times 64 \to 384 \times 64 \times 64.
\]
Under the default \texttt{attention\_resolution = 32}, this stage still does not receive attention.
\paragraph{Block 2.}
The second ResBlock keeps the width:
\[
384 \times 64 \times 64 \to 384 \times 64 \times 64.
\]
\paragraph{Downsample.}
The stage ends with
\[
384 \times 64 \times 64 \to 384 \times 32 \times 32.
\]

\subsection{Encoder stage 3: $32 \times 32$}
This is the first encoder stage where attention is active for the default threshold.
The incoming feature map is
\[
384 \times 32 \times 32.
\]
The stage embedding is
\[
\ell_3 = \mathrm{joint} + W_3(c_{\mathrm{emb}}).
\]
\paragraph{Block 1.}
The first ResBlock expands channels:
\[
384 \times 32 \times 32 \to 512 \times 32 \times 32.
\]
Immediately after the block, the model inserts:
\begin{itemize}
  \item a self-attention block, and
  \item a cross-attention block because \texttt{cross\_attn\_cond = true}.
\end{itemize}
The cross-attention uses
\[
\text{query} = \text{feature map}, \qquad \text{context} = c_{\mathrm{emb}}.
\]
\paragraph{Block 2.}
The second ResBlock keeps the width:
\[
512 \times 32 \times 32 \to 512 \times 32 \times 32.
\]
Self-attention and cross-attention are appended again.
\paragraph{Downsample.}
The stage ends with
\[
512 \times 32 \times 32 \to 512 \times 16 \times 16.
\]

\subsection{Bottleneck: $16 \times 16$}
The bottleneck is the most compressed representation:
\[
512 \times 16 \times 16.
\]
The bottleneck embedding is
\[
\ell_{\mathrm{bot}} = \mathrm{joint} + W_{\mathrm{bot}}(c_{\mathrm{emb}}).
\]
\paragraph{Block 1.}
The first bottleneck ResBlock keeps the width:
\[
512 \times 16 \times 16 \to 512 \times 16 \times 16.
\]
\paragraph{Attention and ChannelSE.}
The bottleneck then applies:
\begin{itemize}
  \item self-attention,
  \item cross-attention with \texttt{cross\_ctx = c\_emb},
  \item ChannelSE, which reweights channels but does not directly consume the condition,
\end{itemize}
followed by one more ResBlock:
\[
512 \times 16 \times 16 \to 512 \times 16 \times 16.
\]

\subsection{Decoder stage 0: $16 \times 16 \to 32 \times 32$}
The decoder starts by upsampling the bottleneck feature:
\[
512 \times 16 \times 16 \to 512 \times 32 \times 32.
\]
The corresponding encoder skip from stage 3 is also
\[
512 \times 32 \times 32.
\]
These are concatenated to obtain
\[
1024 \times 32 \times 32.
\]
The stage embedding is
\[
\ell^{\mathrm{dec}}_0 = \mathrm{joint} + W^{\mathrm{dec}}_0(c_{\mathrm{emb}}).
\]
The decoder stage contains \texttt{num\_res\_blocks + 1 = 3} ResBlocks:
\begin{itemize}
  \item the first block reduces width from 1024 to 512,
  \item the second and third blocks keep 512 channels.
\end{itemize}
Because this stage is at the smallest decoder resolution, attention is active here for the default threshold.

\subsection{Decoder stage 1: $32 \times 32 \to 64 \times 64$}
The next decoder stage upsamples:
\[
512 \times 32 \times 32 \to 512 \times 64 \times 64.
\]
It then concatenates the encoder skip
\[
384 \times 64 \times 64
\]
to form
\[
896 \times 64 \times 64.
\]
The stage embedding is
\[
\ell^{\mathrm{dec}}_1 = \mathrm{joint} + W^{\mathrm{dec}}_1(c_{\mathrm{emb}}).
\]
The first ResBlock reduces width to 384, and the next two blocks keep 384 channels.
Under the default bookkeeping threshold, this decoder stage still receives attention.

\subsection{Decoder stage 2: $64 \times 64 \to 128 \times 128$}
The feature map is now
\[
384 \times 64 \times 64.
\]
It is upsampled to
\[
384 \times 128 \times 128,
\]
then concatenated with the encoder skip
\[
256 \times 128 \times 128,
\]
giving
\[
640 \times 128 \times 128.
\]
The stage embedding is
\[
\ell^{\mathrm{dec}}_2 = \mathrm{joint} + W^{\mathrm{dec}}_2(c_{\mathrm{emb}}).
\]
The first ResBlock reduces width to 256 and the next two keep 256.
At this resolution, attention is no longer active for the default threshold of 32.

\subsection{Decoder stage 3: $128 \times 128 \to 256 \times 256$}
The final decoder stage upsamples
\[
256 \times 128 \times 128 \to 256 \times 256 \times 256,
\]
then concatenates the first encoder skip
\[
128 \times 256 \times 256,
\]
to obtain
\[
384 \times 256 \times 256.
\]
The stage embedding is
\[
\ell^{\mathrm{dec}}_3 = \mathrm{joint} + W^{\mathrm{dec}}_3(c_{\mathrm{emb}}).
\]
The first ResBlock reduces the width to 128, and the next two keep 128.
No attention is used at this resolution.

\subsection{Output head}
The final decoder feature is
\[
128 \times 256 \times 256.
\]
It is passed through
\[
\mathrm{out} = \mathrm{Conv}_{3 \times 3}(\mathrm{SiLU}(\mathrm{GN}(h))),
\]
which returns the final 3-channel tensor
\[
3 \times 256 \times 256.
\]

\subsection{Why this walkthrough matters}
This is the exact route the tensor follows in the code:
\begin{enumerate}
  \item encode local structure at full resolution,
  \item compress toward a global bottleneck,
  \item apply attention and conditioning at low resolution,
  \item decode back using skip connections,
  \item reconstruct the original 3-channel map.
\end{enumerate}

The key conditioning design points are:
\begin{itemize}
  \item the root \texttt{joint} embedding always feeds AdaGN,
  \item stage-specific projections let each scale interpret the same physics differently,
  \item cross-attention injects the condition again at low resolution,
  \item ChannelSE reweights channels but does not consume the condition directly.
\end{itemize}

\end{document}
